\documentclass[10pt,aspectratio=169]{beamer}

% --- Minimal theme ---------------------------------------------------------
\usetheme{default}
\usecolortheme{dove}
\usefonttheme{professionalfonts}
\setbeamertemplate{navigation symbols}{}
\setbeamertemplate{footline}[frame number]
\setbeamertemplate{itemize items}[circle]
\setbeamertemplate{section in toc}[sections numbered]
\setbeamerfont{frametitle}{size=\large,series=\bfseries}
\setbeamercolor{frametitle}{fg=black}
\setbeamercolor{title}{fg=black}
\setbeamercolor{structure}{fg=black}

\usepackage{booktabs}
\usepackage{array}
\usepackage{graphicx}
\graphicspath{{../figs/}}
\usepackage{amsmath}

\title{Temperature-dependent measurement-error variance}
\subtitle{Does $\sigma_k(T)$ depend on temperature? Step 1 --- diagnostic.}
\author{}
\date{\today}

\begin{document}

\begin{frame}\titlepage\end{frame}

\begin{frame}{Motivation}
  \begin{itemize}
    \item Our factor model in \texttt{synth\_compare.R} assumed per-country
          but temperature-homoskedastic measurement error:
          $\mathrm{Var}(\varepsilon_k(c,t)) = \sigma_k^2(c)$.
    \item But two earlier findings hint this may be wrong:
      \begin{itemize}
        \item[(a)] Pairwise bias-vs-level loess (temp-compare deck) is clearly
                   non-flat in $T$.
        \item[(b)] The ERA5 vs MERRA2 ME gap (distance deck) is concentrated
                   in the top temperature decile.
      \end{itemize}
    \item \textbf{Step 1 question}: after fitting the constant-$\sigma$
          factor model, do its residuals show temperature-dependent
          variance?
  \end{itemize}
\end{frame}

\begin{frame}{What we compute}
  From the constant-$\sigma$ EM fit (per-country $\hat F, \hat\lambda_k$),
  take residuals:
  \[
    \hat\varepsilon_k(c, t) = \tilde T_k(c, t) - \hat\lambda_k(c)\,\hat F(c, t).
  \]
  Bin observations into deciles along two axes, pooled over all countries:
  \begin{description}
    \item[Absolute-$T$ (climate axis)] decile of ensemble-mean $T_{\text{bar}}(c,t)$
          --- separates cold from hot \emph{country-years} overall.
    \item[Within-country deviation (weather axis)] decile of
          $T_{\text{bar}}(c,t) - \bar T_{\text{bar}}(c)$ --- separates
          warm from cool \emph{years for a given country}. This is the
          axis that identifies the panel $\hat\beta_2$.
  \end{description}
  Then $\hat\sigma_k(d) = \sqrt{\mathrm{mean}_{(c,t)\in d}\,\hat\varepsilon_k(c,t)^2}$.
\end{frame}

\begin{frame}{Result: $\hat\sigma_k$ by decile, both axes}
  \centering
  \includegraphics[width=\linewidth]{hf_diag_sigma_decile.pdf}

  \vspace{0.2em}
  \footnotesize
  Left: decile of absolute temperature. Right: decile of within-country
  deviation. The two pictures tell different stories.
\end{frame}

\begin{frame}{Climate-axis heteroskedasticity is real and dataset-specific}
  \centering
  \small
  Ratio $\hat\sigma_k(d_{10}) / \hat\sigma_k(d_1)$ on the \textbf{absolute-T axis}:

  \vspace{0.4em}
  \begin{tabular}{l r l}
    \toprule
    Dataset   & Ratio & Direction \\
    \midrule
    MERRA2    & 1.58 & hot $\gg$ cold \\
    UDel      & 1.33 & hot $>$ cold \\
    ERA5-025  & 1.31 & hot $>$ cold \\
    JRA-3Q    & 0.97 & $\approx$ flat \\
    GMFD      & 0.71 & cold $>$ hot \\
    \bottomrule
  \end{tabular}

  \vspace{0.6em}
  \footnotesize
  MERRA2 is $\sim\!60\%$ noisier in the hottest countries; GMFD is actually
  \emph{cleaner} there (it's the cold tail, including high-latitude and
  2010-affected regions, where GMFD has issues). JRA-3Q is nearly flat.
  So the five datasets disagree not just on the level of noise but on its
  direction with $T$.
\end{frame}

\begin{frame}{Weather-axis heteroskedasticity is nearly flat}
  \centering
  \small
  Ratio on the \textbf{within-country-deviation axis}:

  \vspace{0.4em}
  \begin{tabular}{l r}
    \toprule
    Dataset   & Ratio $d_{10}/d_1$ \\
    \midrule
    GMFD      & 1.15 \\
    ERA5-025  & 1.07 \\
    MERRA2    & 1.03 \\
    UDel      & 0.99 \\
    JRA-3Q    & 0.94 \\
    \bottomrule
  \end{tabular}

  \vspace{0.6em}
  \footnotesize
  All ratios within $\pm 15\%$ of one. In \emph{years} that happen to be
  hotter or colder than usual for a country, noise looks similar. The real
  heteroskedasticity is between \emph{countries}, not between years.
\end{frame}

\begin{frame}{Why climate-axis heteroskedasticity is the right one for $\hat\beta_2$}
  Decompose $T_{ct} = \bar T_c + \tilde u_{ct}$. After partialling out the
  country FE, the quadratic regressor is
  \[
    \widetilde{T^2}_{ct} \;\approx\;
      \underbrace{2\,\bar T_c\,\tilde u_{ct}}_{\text{dominant: climate}\times\text{weather}}
      + \underbrace{\tilde u_{ct}^{\,2}}_{\text{sub-dominant}}.
  \]
  So $\hat\beta_2$ is identified from variation in
  $\bar T_c \cdot \tilde u_{ct}$. Under $T_{\text{obs}} = T^* + \eta$, the
  measurement error in this regressor is $2\,\bar T_c\,\eta_{ct}$, with
  conditional variance
  \[
    \mathrm{Var}\bigl(\widetilde{T^2}_{\text{noise}} \,\big|\, c, t\bigr)
    \;\approx\; 4\,\bar T_c^{\,2}\,\sigma_\eta^2(c).
  \]
  Country climate $\bar T_c$ \emph{multiplies} the noise, so
  heteroskedasticity of $\sigma_\eta$ across country climates directly
  scales attenuation of $\hat\beta_2$.
\end{frame}

\begin{frame}{Quantifying the effect: MERRA2 example}
  MERRA2: $\sigma_\eta$ grows by $1.58\times$ from cold to hot countries
  ($d_1 \to d_{10}$ on absolute-$T$ axis). Country climates span roughly
  $\bar T = 5$°C to $\bar T = 28$°C.

  \vspace{0.4em}
  The measurement-error variance in $\widetilde{T^2}$ at a hot country vs a
  cold country:
  \[
    \frac{\mathrm{Var}\bigl(\widetilde{T^2}_{\text{noise}} | \bar T = 28\bigr)}
         {\mathrm{Var}\bigl(\widetilde{T^2}_{\text{noise}} | \bar T = 5\bigr)}
    \;\approx\;
    \underbrace{\left(\frac{28}{5}\right)^2}_{=31\times} \cdot
    \underbrace{\left(\frac{\sigma_\eta(28)}{\sigma_\eta(5)}\right)^2}_{=2.5\times}
    \;\approx\; 78\times.
  \]

  \vspace{0.4em}
  \footnotesize
  The $\bar T^2$ factor alone means attenuation is already a hot-country
  phenomenon. MERRA2's heteroskedasticity makes it $2.5\times$ worse. For
  GMFD, heteroskedasticity runs the other direction (cold countries
  noisier), so attenuation is relatively worse in the cold tail.
\end{frame}

\begin{frame}{Implications: this matters for both synthesis \emph{and} inference}
  \begin{enumerate}
    \item \textbf{Differential attenuation explains part of the hot-tail ME gap.}
          MERRA2 gets more attenuated at hot $T$ than GMFD does; that
          pushes its $\hat\beta_2$-implied ME at $T = 25$°C toward zero
          faster. Consistent with the ranking in the distance deck.
    \item \textbf{T-dependent synthesis is the right object}. Combine the
          five datasets with weights that depend on $\bar T_c$, not
          on weather deviation. At hot countries use JRA-3Q (flat) and
          GMFD (cleaner at hot); at cold countries use ERA5 and JRA-3Q.
    \item \textbf{Attenuation correction on $\hat\beta_2$ is non-trivial}.
          The $\bar T^2$-weighted average noise variance is the relevant
          scalar. For MERRA2 alone it's about $\sim 2.5\times$ the
          unweighted average.
    \item \textbf{Proceed to Step 2}: fit the heteroskedastic factor model
          with $\sigma_k$ indexed by absolute-$T$ decile, pooled across
          countries. Iterative GLS.
  \end{enumerate}
\end{frame}

% ===========================================================================
\section{Step 2 — $\sigma_k(d)$ and inverse-variance weights}

\begin{frame}{Methodological note: iterative GLS is degenerate}
  The textbook Step 2 would fit the heteroskedastic factor model
  $R_k(c,t) = \lambda_k(c)F(c,t) + \varepsilon_k$,
  $\varepsilon_k \sim \mathcal N(0, \sigma_k^2(d(c,t)))$
  by alternating GLS updates of $F$, $\lambda$, and $\sigma_k(d)$.

  \vspace{0.4em}
  \textbf{It doesn't work here.} The cleanest dataset (ERA5, Step-1
  $\sigma \approx 0.065$°C) gets the largest GLS weight $1/\sigma^2$, which
  pulls $F$ toward ERA5 in every $(c,t)$, which shrinks ERA5's residuals,
  which shrinks $\sigma_{\text{ERA5}}$ further. Even with an absolute σ-floor
  of 0.05°C, the fit rests on the floor for ERA5 in every decile. No useful
  heteroskedasticity estimate comes out.

  \vspace{0.4em}
  \textbf{Fix}: use the Step-1 descriptive $\hat\sigma_k(d)$ --- residual
  SDs from the constant-$\sigma$ factor fit. These are biased relative to
  a heteroskedastic MLE but they are unbiased estimates of residual SD per
  (dataset, decile). For constructing inverse-variance synthesis weights,
  that is what we need.
\end{frame}

\begin{frame}{Residual SD per dataset and absolute-$T$ decile}
  \centering
  \includegraphics[width=0.9\linewidth]{hf_sigma_curves.pdf}

  \vspace{0.2em}
  \footnotesize
  ERA5 and JRA-3Q are tightly bundled around 0.07--0.10°C across all deciles.
  MERRA2 rises from 0.14 (cold) to 0.22 (hot). UDel is middle-tier. GMFD is
  noisiest overall and inverted (0.27 cold, 0.19 hot). Five distinct
  noise-vs-climate signatures.
\end{frame}

\begin{frame}{Implied inverse-variance weights $w_k(d)$}
  \centering
  \includegraphics[width=0.95\linewidth]{hf_weights_heatmap.pdf}

  \vspace{0.2em}
  \footnotesize
  Rows sum to 1. The weights answer: if you trust the $\sigma_k(d)$
  estimates and want a minimum-variance combination at temperature decile
  $d$, what mix do you use?
\end{frame}

\begin{frame}{What the weight table tells us}
  \begin{itemize}
    \item \textbf{ERA5} gets 26--51\% weight across deciles --- usually the
          largest single weight. Peaks around decile 9.
    \item \textbf{JRA-3Q} gets 23--44\% --- nearly tied with ERA5, biggest
          in the cold tail (d2: 44\%) and hottest decile (d10: 42\%).
    \item \textbf{MERRA2} is cut from 13\% (cold, d6) to 5.6\% (hot, d10) ---
          downweighted exactly where it's noisiest.
    \item \textbf{GMFD} gets just 2.3\% in d1 (its noisiest decile); 9--10\%
          elsewhere.
    \item \textbf{UDel} is steady at 10--20\%.
  \end{itemize}
  \vspace{0.4em}
  ERA5 and JRA-3Q together carry $\sim 70\%$ of the weight at every
  temperature level --- they dominate, but not overwhelmingly, and the mix
  shifts meaningfully across the T range. That's the hook for Step 3.
\end{frame}

% ===========================================================================
\section{Step 3 — T-dependent MIN-VAR synthesis}

\begin{frame}{Construction of MIN-VAR-T}
  Using the absolute-$T$ decile weights from Step 2, build
  \[
    \hat T_{\text{MIN-VAR-T}}(c, t) \;=\; \sum_{k=1}^{5} w_k\bigl(d(c,t)\bigr) \, T_k(c, t),
  \]
  where $d(c,t)$ is the decile of ensemble-mean $T$ at observation $(c,t)$
  and $w_k(d) \propto 1/\hat\sigma_k^2(d)$, row-normalised to sum to one.
  \vspace{0.3em}

  Key distinction from the constant-$\sigma$ MIN-VAR in
  \texttt{synth\_compare.R}:
  \begin{itemize}
    \item Constant-$\sigma$ uses per-country $\hat\sigma_k^2(c)$ --- 31
          years per country, noisy.
    \item MIN-VAR-T uses pooled $\hat\sigma_k^2(d)$ --- $\sim\!460$ obs per
          decile, well-estimated. And it varies with $T$, which is the
          right axis of heteroskedasticity.
  \end{itemize}
\end{frame}

\begin{frame}{Headline: MIN-VAR-T ties the best and is more stable}
  \centering
  \small
  \begin{tabular}{l r r r r r}
    \toprule
    Method    & $\hat\beta_2\!\times\!10^4$ & SE$\!\times\!10^4$ & $|t|$ & ME(25) & CI width \\
    \midrule
    \textbf{MIN-VAR-T} & $-6.28$ & 1.63 & \textbf{3.86} & $-0.0154$ & \textbf{0.0166} \\
    FACTOR             & $-6.28$ & 1.62 & 3.86          & $-0.0151$ & 0.0171 \\
    MEAN               & $-6.40$ & 1.68 & 3.82          & $-0.0149$ & 0.0170 \\
    ERA5-025 (raw)     & $-5.51$ & 1.46 & 3.79          & $-0.0127$ & 0.0147 \\
    MIN-VAR (const)    & $-4.95$ & 1.48 & 3.34          & $-0.0107$ & 0.0158 \\
    MERRA2 (raw)       & $-4.18$ & 1.41 & 2.97          & $-0.0073$ & 0.0138 \\
    \bottomrule
  \end{tabular}

  \vspace{0.6em}
  \footnotesize
  MIN-VAR-T matches FACTOR on t-stat, beats it on CI width, and --- as the
  next slide shows --- on jackknife stability. The constant-$\sigma$
  MIN-VAR from our previous deck is the \emph{worst} synthesis after
  attenuation; switching to pooled $\sigma_k(d)$ fixes this.
\end{frame}

\begin{frame}{Coefficient plot: MIN-VAR-T joins the bundle}
  \centering
  \includegraphics[width=0.85\linewidth]{hf_forest_update.pdf}

  \vspace{0.2em}
  \footnotesize
  MIN-VAR-T (highlighted, larger marker) sits at the deattenuated end
  of the bundle, essentially coincident with FACTOR and MEAN.
\end{frame}

\begin{frame}{Distribution of applied weights per observation}
  \centering
  \includegraphics[width=\linewidth]{hf_weight_distribution.pdf}

  \vspace{0.2em}
  \footnotesize
  Each panel shows, across all 4{,}592 observations, the distribution of
  weight assigned to that dataset by MIN-VAR-T. ERA5 sits at 25--50\% mass;
  GMFD concentrates below 10\%; MERRA2 has a visible bimodal split between
  cold-tail observations ($\sim 10$\%) and hot-tail ones ($\sim 6$\%).
\end{frame}

\begin{frame}{Jackknife stability: leave-one-dataset-out}
  \centering
  \small
  ME(25°C) after dropping each dataset and rebuilding MIN-VAR-T from scratch:

  \vspace{0.4em}
  \begin{tabular}{l r r}
    \toprule
    Dataset dropped & ME(25°C) & $\hat\beta_2 \times 10^4$ \\
    \midrule
    ERA5-025 & $-0.0152$ & $-6.17$ \\
    GMFD     & $-0.0149$ & $-6.08$ \\
    JRA-3Q   & $-0.0142$ & $-6.06$ \\
    MERRA2   & $-0.0148$ & $-6.09$ \\
    UDel     & $-0.0151$ & $-6.19$ \\
    \midrule
    \textbf{Range}   & \textbf{0.00097} & --- \\
    \bottomrule
  \end{tabular}

  \vspace{0.6em}
  \footnotesize
  Range of 0.001 vs 0.0015 for FACTOR and 0.0016 for MEAN (from
  \texttt{synth\_compare} deck). MIN-VAR-T is the most jackknife-stable
  synthesis we have.
\end{frame}

% ===========================================================================
\section{Step 4 — EIV-corrected damage function}

\begin{frame}{Classical EIV correction with T-dependent noise}
  Apply the standard errors-in-variables formula to the residualised
  regressors $\widetilde T$ and $\widetilde{T^2}$:
  \[
    \hat\beta^{\text{corr}} = \hat\beta_{\text{obs}} / \kappa, \qquad
    \kappa = \frac{\mathrm{Var}(\widetilde X_{\text{obs}}) - \mathrm{Var}(\widetilde X_{\text{noise}})}
                   {\mathrm{Var}(\widetilde X_{\text{obs}})}.
  \]
  For MIN-VAR-T, per-observation noise variance is
  $\bar\sigma^2(d(c,t)) = \bigl(\sum_k 1/\hat\sigma_k^2(d)\bigr)^{-1}$.
  The measurement-error variance in $\widetilde{T^2}$ is dominated by the
  cross term:
  \[
    \mathrm{Var}(\widetilde{T^2}_{\text{noise}}) \approx
      4\,\mathbb E\bigl[\,\bar T_c^{\,2}\,\bar\sigma^2(d(c,t))\,\bigr].
  \]
  Apply $\kappa_1$ to $\hat\beta_1$, $\kappa_2$ to $\hat\beta_2$; SEs scale
  by $1/\kappa$.
\end{frame}

\begin{frame}{Local identification strength $\kappa(d)$}
  \centering
  \includegraphics[width=0.85\linewidth]{hf_attenuation_factor.pdf}

  \vspace{0.2em}
  \footnotesize
  For MIN-VAR-T, $\kappa(d) \ge 0.98$ at every decile. The synthesis
  collapsed the effective measurement-error variance to $\sim 0.04$--$0.06$°C
  per observation, so the remaining attenuation is tiny.
\end{frame}

\begin{frame}{Corrected ME curve under MIN-VAR-T: a small shift}
  \centering
  \includegraphics[width=0.85\linewidth]{hf_me_corrected.pdf}

  \vspace{0.2em}
  \small
  \begin{tabular}{l r r r}
    \toprule
                  & $\hat\beta_2 \times 10^4$ & ME(25°C) & amplification \\
    \midrule
    MIN-VAR-T obs.      & $-6.28$ & $-0.0154$ & --- \\
    EIV-corrected       & $-6.53$ & $-0.0164$ & $\times 1.06$ \\
    \bottomrule
  \end{tabular}

  \vspace{0.4em}
  \footnotesize
  $\kappa_1 = 0.983$, $\kappa_2 = 0.962$. Only 4\% of the $\widetilde{T^2}$
  variance is noise. Correction is real but modest.
\end{frame}

\begin{frame}{But for raw datasets, attenuation is substantial}
  \centering
  \small
  \begin{tabular}{l r r r r}
    \toprule
    Method       & $\kappa_2$ & $\hat\beta_2$ obs & $\hat\beta_2$ corr & amplification \\
    \midrule
    MIN-VAR-T    & 0.962 & $-6.28\!\times\!10^{-4}$ & $-6.53\!\times\!10^{-4}$ & $\times 1.04$ \\
    ERA5-025     & 0.907 & $-5.51\!\times\!10^{-4}$ & $-6.08\!\times\!10^{-4}$ & $\times 1.10$ \\
    JRA-3Q       & 0.892 & $-6.22\!\times\!10^{-4}$ & $-6.97\!\times\!10^{-4}$ & $\times 1.12$ \\
    UDel         & 0.771 & $-5.17\!\times\!10^{-4}$ & $-6.70\!\times\!10^{-4}$ & $\times 1.30$ \\
    MERRA2       & 0.679 & $-4.18\!\times\!10^{-4}$ & $-6.16\!\times\!10^{-4}$ & $\times 1.47$ \\
    GMFD         & 0.551 & $-6.43\!\times\!10^{-4}$ & $-11.7\!\times\!10^{-4}$ & $\times 1.82$ \\
    \bottomrule
  \end{tabular}

  \vspace{0.4em}
  \footnotesize
  MERRA2's $\hat\beta_2$ was attenuated 32\%; GMFD's, 45\%. MIN-VAR-T is
  only 4\% attenuated because the synthesis has already collapsed the
  measurement error.
\end{frame}

\begin{frame}{Headline: correction largely reconciles the datasets}
  \centering
  \includegraphics[width=0.85\linewidth]{hf_me_corrected_compare.pdf}

  \vspace{0.2em}
  \footnotesize
  ME(25°C) before correction: MERRA2 at $-0.007$, JRA-3Q at $-0.016$ ---
  a factor-of-2 spread. After EIV correction the spread narrows
  substantially, with most methods (MIN-VAR-T, MERRA2, ERA5, UDel) landing
  in $[-0.018, -0.014]$. GMFD over-corrects to $-0.028$, probably signalling
  non-classical measurement error in that dataset (consistent with the 2010
  anomaly we flagged in the temp-compare deck).
\end{frame}

% ===========================================================================
\section{Step 5 --- Takeaways}

\begin{frame}{What we learned}
  \begin{enumerate}
    \item \textbf{Heteroskedasticity is a climate-axis phenomenon, not a
          weather-axis one.} Within-country decile ratios $d_{10}/d_1$
          are all within $\pm 15\%$ of 1. Between-country decile ratios
          range from 0.71 (GMFD, inverted) to 1.58 (MERRA2). The five
          products have distinct noise-vs-climate signatures.

    \item \textbf{Climate-axis $\sigma$ matters for $\hat\beta_2$ through
          the $T^2$ cross term.} After country FE,
          $\widetilde{T^2}_{ct} \approx 2\bar T_c \tilde u_{ct}$. The
          measurement-error variance in this regressor is
          $4\bar T_c^2\,\sigma_\eta^2(\bar T_c)$ --- $\bar T_c$ multiplies
          the noise, and heteroskedasticity of $\sigma_\eta$ compounds on top.

    \item \textbf{T-dependent MIN-VAR is the best synthesis.}
          $t(\hat\beta_2) = 3.86$ (tied top); CI width 0.0166 (tightest);
          jackknife range 0.00097 (most stable). ERA5 and JRA-3Q carry
          $\sim\!70\%$ of the weight across all deciles; MERRA2 and GMFD
          get downweighted where each is noisiest.

    \item \textbf{EIV correction collapses the cross-dataset ME spread.}
          Raw $\hat\beta_2$'s are attenuated 10--45\% (MERRA2 was 32\%,
          GMFD 45\%). Post-correction, four of six methods converge to
          ME(25°C) $\in [-0.018, -0.014]$. The ``hot-tail ME gap'' is
          mostly attenuation bias.

    \item \textbf{Synthesis does most of the work of EIV correction.}
          MIN-VAR-T's $\kappa_2 = 0.96$ (only 4\% attenuation) vs MERRA2
          alone at 0.68 (32\%). A good synthesis and a good EIV correction
          are substitutes on this data.
  \end{enumerate}
\end{frame}

\begin{frame}{Caveats and follow-ups}
  \begin{itemize}
    \item \textbf{$\sigma_k(d)$ is descriptive, not MLE.} The iterative GLS
          fit was degenerate (ERA5's $\sigma$ collapsed to the floor). We
          used residuals of the constant-$\sigma$ factor fit. A clean
          heteroskedastic MLE would need stronger regularisation
          (e.g.\ Bayesian prior on $\sigma$, shared loadings).
    \item \textbf{Classical EIV on a quadratic regressor is an approximation.}
          We treat noise in $\widetilde T$ and $\widetilde{T^2}$ as separate
          additive errors and apply scalar $\kappa_1$, $\kappa_2$
          independently. A rigorous correction would use the full
          $2 \times 2$ noise-covariance matrix.
    \item \textbf{GMFD over-corrects.} $\times 1.82$ correction on
          $\hat\beta_2$ pushes the implied ME to $-0.028$, an outlier.
          Suggests non-classical measurement error (the 2010 anomaly is a
          known candidate; a station-coverage transition around then is
          another).
    \item \textbf{Pooled $\sigma_k(d)$ assumes homogeneity across countries
          within decile.} Continent-pooled $\sigma_k$ or parametric
          $\sigma_k(T) = \sigma_{0k}\exp(\gamma_k T)$ is a robustness check.
    \item \textbf{Daily data.} The hot-tail story is fundamentally about
          days above $30$°C; annual means smooth this out. Re-running this
          on daily data is the obvious next frontier.
  \end{itemize}
\end{frame}

\begin{frame}{Bottom line}
  \begin{center}
    \Large
    The ME gap between reanalyses is measurement error.\\[0.5em]
    A T-dependent inverse-variance synthesis plus classical EIV correction\\
    bring four of five datasets into agreement at the hot tail.\\[0.8em]
    \normalsize
    MIN-VAR-T + EIV gives $\hat\beta_2 = -6.53\times10^{-4}$,\\
    ME($25$°C) $= -0.016$, with the tightest CI and the most stable\\
    jackknife profile of any method we have tested.
  \end{center}
\end{frame}

\end{document}
